\SourceLeaf{166}
\begin{center}
{\large NACHLASS UND BRIEFWECHSEL.\par}
\medskip
[ÜBER DIE ERSTEN GRÜNDE DER GEOMETRIE.]
\end{center}

\bigskip
\begin{center}\textsc{Gauss} an \textsc{Bolyai}. Helmstedt, 16. Dezember 1799.\end{center}

\medskip
.\, .\, .\, .\, .\, .\, . Es tut mir sehr leid, dass ich unsere frühere größere Nähe nicht genutzt habe, um mehr von Deinen Arbeiten über die ersten Gründe der Geometrie zu erfahren. Ich hätte mir dadurch gewiss manche vergebliche Mühe erspart und wäre ruhiger geworden, als jemand wie ich es sein kann, solange bei einem solchen Gegenstand noch so viel zu wünschen übrig bleibt. Ich selbst bin in meinen Arbeiten darüber weit vorangekommen (obwohl mir meine anderen, ganz andersartigen Geschäfte wenig Zeit dafür lassen). Der Weg aber, den ich eingeschlagen habe, führt weniger zu dem Ziel, das man sich wünscht und das Du erreicht zu haben versicherst, als vielmehr dazu, die Wahrheit der Geometrie in Zweifel zu ziehen. Zwar bin ich auf manches gestoßen, was bei den meisten schon als Beweis gelten würde, in meinen Augen aber so gut wie NICHTS beweist. Könnte man z.\,B. beweisen, dass ein geradliniges Dreieck möglich ist, dessen Flächeninhalt größer als jede gegebene Fläche ist, so könnte ich die ganze Geometrie völlig streng beweisen. Die meisten würden das wohl als Axiom gelten lassen. Ich nicht. Es wäre ja möglich, dass der Flächeninhalt stets unter (infra) einer gegebenen Grenze bleibt, wie weit man auch die drei Eckpunkte des Dreiecks im Raum voneinander entfernt annimmt. Solche Sätze habe ich mehrere, finde aber keinen davon befriedigend. Mach doch
